% ── 다리 초안: baek_pilon2022 → 우리 엔트로피계수 판독틀(§1 bgbox (ii) / Part T) ──
% 삽입 위치: ch1_sec01_n0n1.tex, 측정원리 bgbox(현행 204–218행) 직후,
%   br_weppner_huggins1977.tex 스니펫과 인접(그 뒤). 자산 주석(현행 219행) 앞.
% 앵커 문장: "...상전이 유형(균질 고용체·이온 질서화·두-상 공존)별로 읽는 해석 지도는 리뷰
%   \cite{baek_pilon2022} 가 정리한다..."(현행 211–214행)
% 컨테이너: srcbox. 우리 기호로 서술. 유도식 아닌 해석-지도 다리(대응표 + 방법 + 가정 차).
% 서지: tier B → load-bearing 판독엔 1차 문헌(reynier2003·allart2018) 병기(bib 규약).

\begin{srcbox}
\textbf{다리 --- 엔트로피계수 부호 판독이 어느 해석 지도인가.} 본문이 $\partial U_\oc/\partial T=\Delta
S(x)/F$ 의 부호$\cdot$모양을 상전이 유형별로 읽는 틀은 Baek--Pilon 의 전위차 엔트로피 측정 해석
지도\cite{baek_pilon2022} 다:
\begin{center}\footnotesize
\begin{tabular}{@{}ll@{}}
\hline
본문 & Baek--Pilon\cite{baek_pilon2022} 대응 \\
\hline
측정 기울기 $\partial U_\oc/\partial T=\Delta S(x)/F$ & entropic potential $\partial U/\partial T$ \\
단상 폭(\S\ref{sec:width}) & 균질 고용체 서명 \\
두-상 broadening(\S\ref{sec:broadening}) & 두-상 공존 서명 \\
Part T 의 SOC 부호 교대 & 유형별 부호$\cdot$모양 서명 \\
\hline
\end{tabular}
\end{center}
\emph{방법 요지} --- Baek--Pilon 은 전위차 엔트로피(entropic potential) 측정 $\partial U/\partial T$ 의
부호$\cdot$모양을 상전이 유형(균질 고용체$\cdot$이온 질서화$\cdot$두-상 공존)별 열역학 서명으로 갈라 읽는 해석
지도를 세운다. \emph{가정 차} --- 이는 측정 해석 리뷰(tier B)라 특정 전극의 미시 모형을 규정하지 않으므로,
본문은 유형 분류만 판독틀로 빌리고 실제 부호$\cdot$크기는 우리 세 분포 합(config$+$vib$+$elec, Part T)과 1차
문헌(\cite{reynier2003,allart2018})으로 확정한다(tier B 인용의 1차 병기 규약).
\end{srcbox}
